\documentclass[conference]{IEEEtran}

\usepackage{amsmath}
\usepackage{booktabs}
\usepackage{cite}
\usepackage{graphicx}
\usepackage{microtype}
\usepackage{siunitx}
\usepackage[hidelinks]{hyperref}

\graphicspath{{figures/}}
\sisetup{detect-all,round-mode=places,round-precision=3}
% Generated by scripts/build_assets.py. Do not edit by hand.
\newcommand{\BreathValidationAUROC}{0.783}
\newcommand{\CoughValidationAUROC}{0.807}
\newcommand{\SpeechValidationAUROC}{0.843}
\newcommand{\BreathTestAUROC}{0.828}
\newcommand{\CoughTestAUROC}{0.862}
\newcommand{\SpeechTestAUROC}{0.888}
\newcommand{\SpeechTestAUPRC}{0.833}
\newcommand{\SpeechTestBalancedAccuracy}{0.809}
\newcommand{\SpeechTestFOne}{0.740}
\newcommand{\FusionValidationAUROC}{0.842}
\newcommand{\FusionValidationAUPRC}{0.800}
\newcommand{\AllModalitiesValidationAUROC}{0.841}
\newcommand{\AllModalitiesValidationAUPRC}{0.804}
\newcommand{\AllModalitiesTestAUROC}{0.890}
\newcommand{\AllModalitiesTestAUPRC}{0.847}
\newcommand{\BreathSpeechValidationAUROC}{0.835}
\newcommand{\BreathSpeechValidationAUPRC}{0.797}
\newcommand{\BreathSpeechTestAUROC}{0.886}
\newcommand{\BreathSpeechTestAUPRC}{0.824}
\newcommand{\CoughBreathValidationAUROC}{0.819}
\newcommand{\CoughBreathValidationAUPRC}{0.775}
\newcommand{\CoughBreathTestAUROC}{0.867}
\newcommand{\CoughBreathTestAUPRC}{0.816}
\newcommand{\FusionTestAUROC}{0.895}
\newcommand{\FusionTestAUPRC}{0.862}
\newcommand{\FusionTestBalancedAccuracy}{0.808}
\newcommand{\FusionTestFOne}{0.729}
\newcommand{\FusionTestSensitivity}{0.843}
\newcommand{\FusionTestSpecificity}{0.774}
\newcommand{\FusionTestN}{314}
\newcommand{\FusionThreshold}{0.325}
\newcommand{\FusionAUCILow}{0.852}
\newcommand{\FusionAUCIHigh}{0.933}
\newcommand{\WeightedTestAUROC}{0.896}
\newcommand{\StackTestAUROC}{0.897}
\newcommand{\CoughFusionWeight}{0.459}
\newcommand{\SpeechFusionWeight}{0.541}
\newcommand{\FusionMinusSpeechAUROC}{0.007}
\newcommand{\FusionMinusSpeechAUPRC}{0.028}
\newcommand{\CandidateFeatureCount}{10{,}140}
\newcommand{\CompareFeatureCount}{6{,}373}
\newcommand{\ISFeatureCount}{1{,}582}
\newcommand{\ProjectFeatureCount}{2{,}177}
\newcommand{\OpenSmileEventFeatureCount}{8}
\newcommand{\SelectedFeatureCount}{800}
\newcommand{\IndexedParticipants}{2{,}746}
\newcommand{\ResolvedParticipants}{2{,}114}
\newcommand{\QualityParticipants}{2{,}088}
\newcommand{\TrainParticipants}{1{,}460}
\newcommand{\ValidationParticipants}{312}
\newcommand{\TestParticipants}{316}
\newcommand{\TrainPositive}{476}
\newcommand{\ValidationPositive}{101}
\newcommand{\TestPositive}{103}
\newcommand{\TrainNegative}{984}
\newcommand{\ValidationNegative}{211}
\newcommand{\TestNegative}{213}
\newcommand{\TrainRecordings}{12{,}685}
\newcommand{\ValidationRecordings}{2{,}702}
\newcommand{\TestRecordings}{2{,}719}


\hypersetup{
  pdftitle={Validation-Guided Cough-Speech Fusion for Participant-Disjoint COVID-19 Classification from Respiratory Audio},
  pdfauthor={Nishant Harkut and Ritu Tiwari},
  pdfkeywords={COVID-19, respiratory audio, multimodal fusion, participant-disjoint evaluation}
}

\title{Validation-Guided Cough--Speech Fusion for Participant-Disjoint COVID-19 Classification from Respiratory Audio}

\author{
\IEEEauthorblockN{Nishant Harkut and Ritu Tiwari}
\IEEEauthorblockA{Department of Information Technology\\
Atal Bihari Vajpayee Indian Institute of Information Technology and Management\\
Gwalior, India}
}

\begin{document}
\maketitle

\begin{abstract}
Crowdsourced respiratory-audio datasets contain several recordings per
participant, enabling audio-only fusion but also creating opportunities for
participant leakage and test-guided model choice. We developed a Coswara
system in fixed participant-disjoint train, validation, and test partitions.
Quality-screened cough, breathing, and speech recordings were represented by
\CandidateFeatureCount{} numeric acoustic candidates. LightGBM feature ranking
was fitted only on training data and retained \SelectedFeatureCount{}
candidates; validation participants selected modality models, the modality
combination, and the operating threshold. The selected speech branch achieved
test AUROC \SpeechTestAUROC{}. Uniformly averaging the selected cough and
speech probabilities yielded AUROC \FusionTestAUROC{} (95\% participant-bootstrap
CI \FusionAUCILow{}--\FusionAUCIHigh{}), AUPRC \FusionTestAUPRC{}, balanced
accuracy \FusionTestBalancedAccuracy{}, and F1 \FusionTestFOne{} on
\FusionTestN{} participants. The AUROC difference from speech alone was
\FusionMinusSpeechAUROC{} and is reported descriptively because paired
uncertainty for the difference was unavailable. The study establishes a
reproducible validation-guided fusion workflow for internal source evaluation;
it does not establish clinical or cross-dataset validity.
\end{abstract}

\begin{IEEEkeywords}
COVID-19, respiratory audio, multimodal fusion, participant-disjoint
evaluation, acoustic features
\end{IEEEkeywords}

\section{Introduction}

Smartphone recordings of cough, breathing, and speech have been investigated
as accessible inputs for COVID-19 risk assessment. A systematic review found
substantial variation across the field in datasets, preprocessing, acoustic
representations, classifiers, and evaluation procedures
\cite{aleixandre2022review}. Public resources also differ in what they collect:
COVID-19 Sounds contains cough, breath, and voice \cite{xia2021covid}; COUGHVID
is a large cough corpus \cite{orlandic2021coughvid}; and Coswara records nine
prompted respiratory and vocal sound tasks alongside participant metadata
\cite{bhattacharya2023coswara}. Consequently, the term \emph{multimodal} can
refer to different sound tasks, learned representations, or the addition of
non-audio variables.

Multiple sound tasks are attractive because infection-related changes need not
appear identically in a forced cough, sustained phonation, and connected speech.
Prior Coswara studies used transfer learning across cough, breath, and speech
\cite{pahar2022transfer}, combined acoustic categories and symptoms
\cite{chetupalli2023multimodal}, or fused waveform- and spectrogram-derived
representations across several recordings \cite{truong2024fair}. These results
support multisound analysis, but they do not imply that every available stream
should be included. An added sound type can increase missingness, dimensionality,
and selection variance; its value should therefore be established on participants
who were not used to fit the modality models.

Split construction and model selection must also be reported as part of the
experimental method. Han \emph{et al.} showed that participant allocation and
demographic composition can substantially alter respiratory-audio estimates
\cite{han2022realistic}. Coppock \emph{et al.} found that matching on recruitment
channel, demographics, and symptoms reduced apparent audio performance
\cite{coppock2024audio}. Deployment and longitudinal studies further document
sensitivity to acquisition setting and calendar period
\cite{grant2022deployment,kim2024longitudinal,ganitidis2025drift}. These concerns
do not prevent internal model development, but they require a clear separation
between fitting, model selection, and final estimation.

We therefore evaluated which audio-only probability-fusion configuration a
fixed participant-disjoint Coswara protocol selects from validation evidence,
and compared its held-out result descriptively with its constituent modalities.
Here \emph{multimodal} means multiple respiratory sound types from the same
participant; symptoms and demographic variables are not model inputs. The
contributions are: 1) quality-controlled processing of cough, breathing, and
speech under participant-isolated partitions; 2) training-only ranking of a
combined openSMILE and signal-summary feature bank; 3) validation-only selection
of modality branches, modality combinations, and the operating threshold; and
4) a participant-level held-out estimate with bootstrap uncertainty and explicit
inferential limits.

\section{Related Work}

\subsection{Datasets and Benchmark Tasks}

Early crowdsourced studies established that cough and breath recordings could
be modeled with measured descriptors or pretrained audio embeddings
\cite{brown2020exploring}. The DiCOVA challenge subsequently provided a common
blind-test benchmark with separate cough and multisound tracks
\cite{muguli2021dicova}. Coswara differs from a cough-only benchmark because one
participant may provide shallow and heavy coughs, shallow and deep breathing,
counting speech, and sustained vowels. This repeated-measures structure enables
participant-aligned fusion, but it also makes recording-level random splitting
inappropriate. Dataset name alone is therefore insufficient for comparison:
the label definition, participant snapshot, selected sound tasks, and analysis
unit must also agree.

\subsection{Representations and Fusion}

Published systems span end-to-end spectrogram networks, pretrained embeddings,
and engineered acoustic descriptors. Laguarta \emph{et al.} used ResNet-based
cough branches \cite{laguarta2020cough}; Pahar \emph{et al.} examined global
smartphone cough recordings and later transfer-learning bottlenecks for cough,
breath, and speech \cite{pahar2021smartphone,pahar2022transfer}; and
CovidCoughNet evaluated convolutional representations across cough, breath, and
voice \cite{celik2023covidcoughnet}. Feature-engineered alternatives remain
competitive in small audio cohorts. Chowdhury \emph{et al.} combined ranked
acoustic descriptors with ensemble decision making \cite{chowdhury2022mcdm},
whereas Avila \emph{et al.} fused a ComParE-SVM and log-mel CNN on the DiCOVA
blind test \cite{avila2021feature}. Hierarchical spectrogram transformers offer
a higher-capacity alternative for individual respiratory sound tasks
\cite{aytekin2024hst}.

Fusion can occur before classification by concatenating or jointly attending to
representations, or after classification by combining modality probabilities
\cite{baltrusaitis2019survey}. Chetupalli \emph{et al.} reported 0.880 AUC from
three acoustic categories and 0.963 after adding symptoms
\cite{chetupalli2023multimodal}. FAIR combined frozen wav2vec and DeiT
representations with attention over seven Coswara recordings and reported
$0.8658\pm0.0115$ AUC \cite{truong2024fair}. The present study uses score-level
fusion so that each sound branch can be inspected, selected, and aggregated at
participant level without learning a new high-capacity representation on the
final test cohort.

\subsection{Evaluation Practice}

The literature contains both participant-separated evaluation and protocols in
which grouping is unclear. Feature selection can be another source of optimism:
Avila \emph{et al.} documented a marked difference between nonnested development
estimates and a blind evaluation \cite{avila2021feature}. Han \emph{et al.} and
Coppock \emph{et al.} further demonstrated that cohort construction changes the
question answered by an audio classifier \cite{han2022realistic,coppock2024audio}.
Following the reporting principles emphasized by TRIPOD+AI
\cite{collins2024tripod}, we state the independent unit, data role, selection
source, threshold source, and uncertainty target. Literature values are used as
protocol context, not as interchangeable leaderboard entries.

\section{Materials and Methods}

\subsection{Cohort, Labels, and Partitioning}

The indexed Coswara release contained 2,746 participants. We retained 2,114
participants with resolved active-positive or healthy/negative labels; records
with unresolved labels were not fitted or evaluated. Quality screening left
2,088 participants with at least one eligible cough, breathing, or speech
recording. A participant could contribute two coughs, two breathing recordings,
two counting recordings, and three sustained vowels. All recordings from a
participant remained in one fixed partition. The initial 70/15/15 assignment
used seed 42 and stratified by label, age band, and gender where cell sizes
permitted, with a label-only fallback for sparse cells. Table~\ref{tab:cohort}
reports the quality-passing analysis cohort. The participant, rather than a
recording or sound segment, was the independent evaluation unit.

\begin{table}[!t]
\caption{Quality-Passing Participant-Disjoint Partitions}
\label{tab:cohort}
\centering
\footnotesize
\setlength{\tabcolsep}{3.3pt}
\begin{tabular}{lrrrr}
\toprule
Partition & Participants & Positive & Negative & Recordings \\
\midrule
Train      & \TrainParticipants & \TrainPositive & \TrainNegative & \TrainRecordings \\
Validation & \ValidationParticipants & \ValidationPositive & \ValidationNegative & \ValidationRecordings \\
Test       & \TestParticipants & \TestPositive & \TestNegative & \TestRecordings \\
\bottomrule
\end{tabular}
\end{table}

\begin{figure*}[!t]
\centering
\includegraphics[width=\textwidth]{study_design.pdf}
\caption{Participant flow and data responsibilities. Participants with
unresolved labels or no eligible respiratory recording were excluded before
analysis. All recordings from a retained participant were assigned together.
Only training participants informed feature ranking, classifier fitting, and
SMOTE; validation participants determined branches, modality combination, and
threshold; test participants supplied the final locked estimate.}
\label{fig:study-design}
\end{figure*}

\subsection{Audio Quality Control and Preprocessing}

Recordings were decoded as mono waveforms and resampled to 16~kHz. Silence was
trimmed at 35~dB below peak energy. A root-mean-square detector used 25-ms
frames and a 10-ms hop; frames above 25\% of the recording's maximum RMS
defined the active interval, expanded by a 120-ms margin. Peak normalization
then preceded centered cropping or zero padding to 4~s for cough and 5~s for
breathing and speech. A recording was eligible when its duration was at least
0.5~s for cough or 1~s for breathing/speech, its silence ratio did not exceed
0.70, and its clipping ratio did not exceed 0.01. These rules were fixed before
the reported fusion selection. No validation or test waveform was augmented.

\subsection{Acoustic Representation and Feature Ranking}

Three complementary feature blocks were merged by recording identifier. The
openSMILE ComParE~2016 and IS10 functionals contributed
\CompareFeatureCount{} and \ISFeatureCount{} descriptors, respectively,
covering energy, spectral shape, cepstral trajectories, and voicing
\cite{eyben2010opensmile,schuller2010is10,schuller2016compare}. A
\ProjectFeatureCount{}-variable signal-summary block described 40 MFCC
trajectories and their first and second differences, 64 log-mel bands, chroma,
spectral contrast, tonnetz, energy, zero crossing, spectral shape, and tempo.
Each OpenSMILE extraction also carried four prefixed active-interval variables
(start, end, duration, and active-audio ratio), contributing
\OpenSmileEventFeatureCount{} additional columns after merging. Together, the
four rows in Table~\ref{tab:features} account for all \CandidateFeatureCount{}
numeric candidates.

\begin{table}[!t]
\caption{Acoustic Candidate Blocks Before Ranking}
\label{tab:features}
\centering
\footnotesize
\setlength{\tabcolsep}{3.5pt}
\begin{tabular}{lrl}
\toprule
Block & Variables & Principal information \strut \\
\midrule
ComParE 2016 & \CompareFeatureCount & Energy, spectrum, cepstra, voice \\
IS10 & \ISFeatureCount & Paralinguistic functionals \\
Signal summaries & \ProjectFeatureCount & MFCC, mel, chroma, timing \\
Active-interval variables & \OpenSmileEventFeatureCount & Start, end, duration, active ratio \\
\midrule
Merged numeric bank & \CandidateFeatureCount & Complementary descriptors \\
\bottomrule
\end{tabular}
\end{table}

Ranking used only training rows. A LightGBM ranker was fitted separately in
each sound modality. If $g_{jm}$ is the gain of feature $j$ in modality $m$,
the common score was
\begin{equation}
s_j=\frac{1}{|\mathcal{M}|}\sum_{m\in\mathcal{M}}
\frac{g_{jm}}{\max_k g_{km}}.
\label{eq:rank}
\end{equation}
The max normalization prevents one modality's gain scale from dominating the
aggregate. The \SelectedFeatureCount{} highest-scoring names were frozen for
validation and test. Feature budgets of 500, 800, and 1,200 were explored
during development; inference in this paper is specific to the 800-feature
configuration. Within each model pipeline, constant variables were removed
and an ANOVA F-score retained 80\% of candidates for boosted models or 60\%
for the radial-basis SVC.

\subsection{Modality Models and Validation-Guided Fusion}

Four classifier families formed the final candidate bank: LightGBM, XGBoost,
CatBoost, and a sigmoid-calibrated radial-basis SVC
\cite{ke2017lightgbm,chen2016xgboost,prokhorenkova2018catboost,cortes1995svm}.
LightGBM used 900 trees, learning rate 0.025, 31 leaves, minimum child size 20,
row and feature subsampling of 0.85, and $L_2$ regularization 2. XGBoost used
800 depth-3 trees, learning rate 0.025, row and feature subsampling of 0.85,
minimum child weight 2, and $L_2$ regularization 2. CatBoost used 900 depth-5
iterations and learning rate 0.025. The SVC used $C=2$, class-balanced fitting,
scaled inputs, and three-fold sigmoid calibration
\cite{niculescu2005calibration}. For the boosted models, SMOTE with three
nearest neighbors was applied inside the training pipeline after feature
filtering; validation and test data were never resampled \cite{chawla2002smote}.
All stochastic components used seed 42.

Recording probabilities were averaged first within participant and modality.
For each modality, validation AUROC, with AUPRC as the tie breaker, selected
either one classifier or the mean of the four validation-ranked classifiers.
Let $p_m(x)$ denote the selected probability for modality $m$ and participant
$x$. The primary fusion rule for a modality set $C$ was the
uniform mean
\begin{equation}
p_C(x)=\frac{1}{|C|}\sum_{m\in C}p_m(x).
\label{eq:fusion}
\end{equation}
Validation AUROC and AUPRC selected among the four combinations of two or three
sound modalities. A validation-AUPRC-weighted mean and a class-balanced
logistic stack fitted to validation probabilities were retained as exploratory
sensitivity analyses; neither was eligible to replace the uniform primary
rule. For weighted fusion, $a_m=\max(\mathrm{AUPRC}_m-0.5,0.01)$ and
$w_m=a_m/\sum_{k\in C}a_k$, using validation AUPRC. The resulting cough and
speech weights were \CoughFusionWeight{} and \SpeechFusionWeight{}. The probability
threshold maximizing validation balanced accuracy (\FusionThreshold{}) was
then frozen for test balanced accuracy, sensitivity, specificity, and F1.

\begin{figure*}[!t]
\centering
\includegraphics[width=\textwidth]{selection_and_results.pdf}
\caption{Validation-guided selection and locked test evaluation. (a) Uniform
modality combinations ranked on the validation participants; the shaded row was
selected. (b) Participant-level test AUROC and AUPRC for the selected modality
branches and primary cough--speech mean. The horizontal segment is the 95\%
participant-bootstrap interval for the primary fusion AUROC; comparable branch
intervals were unavailable and are not implied.}
\label{fig:selection-results}
\end{figure*}

\subsection{Leakage Controls and Reproducibility}

The data-role boundary was enforced at every supervised step. Split assignment
preceded feature ranking; constant removal, ANOVA filtering, scaling,
calibration, and SMOTE were fitted without validation or test rows. Model and
modality selection used validation outcomes only. The test set was queried for
the locked choices and was not fed back into the candidate bank. This order is
important because selecting features or fusion rules on final outcomes would
turn the reported test value into another development estimate.

Missing sound tasks were not replaced by another participant's recording.
Single-modality metrics used all eligible participants for that modality,
whereas a fusion row used the intersection with valid predictions for every
member of its modality set. The resulting sample size is therefore reported for
each row. Split assignments, selected feature names, model-selection records,
thresholds, and aggregate predictions were retained as machine-readable
artifacts so that each reported number can be traced to its data role.

\subsection{Outcomes and Uncertainty}

The primary endpoint was participant-level AUROC; AUPRC addressed the less
frequent positive class. At the frozen threshold, balanced accuracy was
$(\mathrm{sensitivity}+\mathrm{specificity})/2$, and
$F_1=2(\mathrm{precision}\times\mathrm{recall})/
(\mathrm{precision}+\mathrm{recall})$. Percentile limits for the fixed fusion
AUROC came from 1,000 nonparametric resamples of the \FusionTestN{} test
participants. This quantifies participant-sampling uncertainty for a fixed
system and cohort, not refitting variation. Aligned predictions required for a
paired fusion-versus-speech interval were unavailable, so that difference is
descriptive.

\section{Results}

\subsection{Validation Selection}

Validation AUROC selected the four-model probability mean for breathing
(\BreathValidationAUROC{}) and cough (\CoughValidationAUROC{}), and LightGBM
for speech (\SpeechValidationAUROC{}). Among the uniform combinations,
cough plus speech ranked first (AUROC \FusionValidationAUROC{}, AUPRC
\FusionValidationAUPRC{}), narrowly above cough plus breathing plus speech
(AUROC \AllModalitiesValidationAUROC{}, AUPRC
\AllModalitiesValidationAUPRC{}). Breathing plus speech reached
\BreathSpeechValidationAUROC{}/\BreathSpeechValidationAUPRC{}, and cough plus
breathing reached \CoughBreathValidationAUROC{}/
\CoughBreathValidationAUPRC{} (AUROC/AUPRC). This ordering fixed cough plus
speech before the matching test row was retrieved.

The locked test ordering remained consistent with this choice. Cough plus
speech obtained \FusionTestAUROC{}/\FusionTestAUPRC{}, followed by all three
modalities at \AllModalitiesTestAUROC{}/\AllModalitiesTestAUPRC{}, breathing
plus speech at \BreathSpeechTestAUROC{}/\BreathSpeechTestAUPRC{}, and cough
plus breathing at \CoughBreathTestAUROC{}/\CoughBreathTestAUPRC{}
(AUROC/AUPRC). These post-selection rows are an ablation and were not used to
revise the selected combination.

\subsection{Held-Out Participant Performance}

Table~\ref{tab:results} presents the selected branches and fusion sensitivity
analyses. Speech was the strongest selected single modality with AUROC
\SpeechTestAUROC{} and AUPRC \SpeechTestAUPRC{}. The primary cough--speech
uniform mean yielded AUROC \FusionTestAUROC{} (95\% CI
\FusionAUCILow{}--\FusionAUCIHigh{}), AUPRC \FusionTestAUPRC{}, balanced
accuracy \FusionTestBalancedAccuracy{}, and F1 \FusionTestFOne{} on
\FusionTestN{} complete-case participants. Relative to speech, the descriptive
differences were +\FusionMinusSpeechAUROC{} AUROC and
+\FusionMinusSpeechAUPRC{} AUPRC. The weighted mean and logistic stack gave
AUROCs \WeightedTestAUROC{} and \StackTestAUROC{}, respectively, but were
exploratory because their additional parameters depended on validation data.
At the frozen threshold, the primary fusion sensitivity was
\FusionTestSensitivity{} and specificity was \FusionTestSpecificity{}; these
operating-point values were not retuned on the test cohort.

\begin{table}[!t]
\caption{Participant-Level Held-Out Results}
\label{tab:results}
\centering
\footnotesize
\setlength{\tabcolsep}{3.0pt}
\begin{tabular}{lrrrrr}
\toprule
System & $n$ & AUROC & AUPRC & Bal. Acc. & F1 \\
\midrule
Breathing, four-model mean & 312 & \BreathTestAUROC & 0.734 & 0.745 & 0.654 \\
Cough, four-model mean     & 316 & \CoughTestAUROC & 0.804 & 0.785 & 0.705 \\
Speech, LightGBM           & 314 & \SpeechTestAUROC & \SpeechTestAUPRC & \SpeechTestBalancedAccuracy & \SpeechTestFOne \\
\midrule
\textbf{Cough+speech mean} & \textbf{\FusionTestN} & \textbf{\FusionTestAUROC} & \textbf{\FusionTestAUPRC} & \textbf{\FusionTestBalancedAccuracy} & \textbf{\FusionTestFOne} \\
Cough+speech weighted$^*$  & 314 & \WeightedTestAUROC & 0.863 & 0.820 & 0.744 \\
Cough+speech stack$^*$     & 314 & \StackTestAUROC & 0.863 & 0.825 & 0.752 \\
\bottomrule
\end{tabular}
\par\vspace{1pt}\scriptsize $^*$Exploratory sensitivity analysis; not eligible for primary selection.
\end{table}

\subsection{Context within Internal Coswara Studies}

Table~\ref{tab:context} places the result beside the nearest published
internal Coswara estimates. The Coswara data paper combined nine sound tasks
with symptoms; Chetupalli \emph{et al.} fused three acoustic categories; FAIR
used seven sound instances and two pretrained representations. Cohort snapshot,
label filtering, inputs, and partition construction differ, so the values are
descriptive protocol context rather than a statistical ranking.

\begin{table*}[!t]
\caption{Descriptive Context from Internal Coswara Evaluations}
\label{tab:context}
\centering
\footnotesize
\setlength{\tabcolsep}{4.5pt}
\begin{tabular}{p{2.8cm}p{4.0cm}p{7.7cm}r}
\toprule
Study & Inputs & Evaluation design & AUROC \\
\midrule
Coswara \cite{bhattacharya2023coswara} & Nine sound tasks and symptoms & Random 70/15/15 subject sets & 0.915 \\
Chetupalli \emph{et al.} \cite{chetupalli2023multimodal} & Three acoustic categories & Non-overlapping 80/20 subject sets & 0.880 \\
FAIR \cite{truong2024fair} & Seven sound instances and two representations & Fixed test set; repeated five-fold development evaluation & $0.866\pm0.012$ \\
This work & Cough and speech probabilities & Participant-disjoint 70/15/15 development and test roles & \FusionTestAUROC \\
\bottomrule
\end{tabular}
\end{table*}

\section{Discussion}

Validation selected cough plus speech rather than all three available sound
types. On the held-out participants, uniform fusion was numerically above speech
in AUROC and AUPRC, but the AUROC difference was only
\FusionMinusSpeechAUROC{}; balanced accuracy was essentially unchanged and F1
was lower. Without a paired interval, the evidence does not establish that
fusion is superior to speech. Evaluating every sound combination under the same
validation rule makes the contribution of each additional stream explicit.

The score-level design also makes the fusion weights interpretable. The primary
system assigns equal weight to cough and speech, while the exploratory
validation-AUPRC weighting changes them only to \CoughFusionWeight{} and
\SpeechFusionWeight{}. Its test AUROC differs from the uniform rule by less
than 0.002. The logistic stack is similarly close. Thus, within this cohort,
the choice among these three probability-fusion rules matters less than the
choice of modalities and the separation of participant roles.

The selected combination is consistent with, but does not prove, complementary
information across cough and speech. Cough can emphasize transient respiratory
events, whereas counting and sustained phonation provide longer voiced segments
from which cepstral and voicing statistics can be estimated. Speech was the
strongest individual branch on the locked test set, and adding cough increased
AUROC by only \FusionMinusSpeechAUROC{}. Breathing did not improve the
validation ranking when added to cough and speech. These observations justify
the selected engineering configuration; they are not evidence that any feature
is a disease-specific biomarker.

The result is numerically comparable with internal Coswara estimates from
Chetupalli \emph{et al.} and FAIR, but each study answers a different question.
Chetupalli \emph{et al.} evaluated acoustic-category fusion and separately
showed the effect of adding symptoms \cite{chetupalli2023multimodal}; FAIR used
learned waveform and image representations with attention across seven
recordings \cite{truong2024fair}. Our result instead isolates audio-only
score-level fusion under three explicit participant roles. Accordingly,
Table~\ref{tab:context} supports methodological comparison, not a claim of
statistical superiority or performance ranking.

Three safeguards strengthen the internal estimate. First, participant grouping
prevents recordings from one person appearing in multiple partitions. Second,
feature ranking and imbalance correction use training participants only, while
the test cohort does not select a branch, modality set, fusion rule, or
threshold. Third, recording-level predictions are pooled before metrics are
computed, aligning the analysis unit with the screening subject. The bootstrap
interval then exposes the sampling precision of the primary estimate rather
than presenting a point estimate alone.

Six limitations should be considered. First, this is a retrospective secondary
analysis of one crowdsourced dataset, and the COVID labels were not independently
adjudicated by the present authors. Second, validation data served multiple
selection roles: modality model, modality combination, and threshold. The
validation-first reporting rule was reconstructed after broader development
rather than preregistered, and the 800-feature budget was one of three explored
budgets. A nested outer evaluation would estimate the entire selection procedure
more conservatively. Third, repeated recordings were fitted as separate training
rows before participant-level probability averaging; event timing and
active-audio descriptors may encode recording behavior. Fourth, the fusion
analysis is complete-case and includes 314 test participants. Fifth, the
bootstrap interval covers a fixed model and cohort, while the fusion-minus-speech
difference has no paired interval. Finally, this paper intentionally evaluates
source-domain model development only. Calendar drift, metadata confounding,
calibration, and independent-dataset transport require distinct estimands and
are reserved for a separate, non-overlapping reliability study.

\section{Conclusion}

Training-only acoustic ranking and validation-guided modality selection chose a
uniform cough--speech probability mean. It achieved participant-level test
AUROC \FusionTestAUROC{} compared with \SpeechTestAUROC{} for the selected
speech branch. The small descriptive margin does not support a superiority
claim, but the workflow provides an auditable way to develop multisound models
without participant overlap or test-ranked fusion. Claims beyond internal
Coswara performance require separate temporal and external evaluation.

\section*{Ethics and Data Availability}

This work is a secondary analysis of the public Coswara resource; no new
participants were recruited. Consent, ethics approval, and collection procedures
are documented in the source publication \cite{bhattacharya2023coswara}. Raw
recordings are not redistributed. Analysis code, frozen configuration records,
and aggregate evidence tables will accompany the archived project release.

\emergencystretch=1em
\IEEEtriggeratref{19}
\bibliographystyle{IEEEtran}
\bibliography{references}

\end{document}
